\chapter*{\centering{Tóm tắt đồ án}}
\addcontentsline{toc}{chapter}{Tóm tắt đồ án}

Đồ án này trình bày quá trình thiết kế và triển khai một hệ thống tính toán phân bố
có tính sẵn sàng cao cho ứng dụng microservices trên nền tảng Amazon Elastic Kubernetes Service (EKS).
Mục tiêu của đồ án là xây dựng một kiến trúc hệ thống phân tán có khả năng chịu lỗi,
tự động mở rộng và duy trì dịch vụ ổn định trong các điều kiện tải và sự cố hạ tầng khác nhau.

Trong phạm vi nghiên cứu, nhóm triển khai cụm Kubernetes trên Amazon EKS
theo mô hình đa vùng sẵn sàng (Multi-AZ),
kết hợp các cơ chế điều phối container, cân bằng tải và quản lý vòng đời ứng dụng.
Hệ thống được thiết kế nhằm đảm bảo các thành phần ứng dụng có thể tiếp tục hoạt động
ngay cả khi xảy ra lỗi tại một hoặc nhiều node trong cụm.

Bên cạnh hạ tầng chính, đồ án xây dựng hệ thống giám sát và quan sát
(\textit{Observability}) cho cụm Kubernetes,
bao gồm giám sát tài nguyên và trạng thái hệ thống bằng Prometheus,
trực quan hóa dữ liệu bằng Grafana,
và thu thập, phân tích log tập trung thông qua Grafana Loki.
Hệ thống này giúp theo dõi hiệu năng, phát hiện bất thường
và hỗ trợ phân tích sự cố trong quá trình vận hành.

Kết quả của đồ án cho thấy kiến trúc được đề xuất đáp ứng tốt các yêu cầu
về tính sẵn sàng, khả năng mở rộng và khả năng giám sát,
đồng thời phù hợp với các nguyên lý của hệ thống tính toán phân bố
và kiến trúc microservices hiện đại.
